\begin{figure}[htbp]
\centering
\begin{tikzpicture}

% Data points: (0, 1), (1, 2.5), (2, 1.8), (3, 3.2), (4, 2.9), (5, 4.1)
% These will be used across all plots

% ========== Data Points Only ==========
\begin{axis}[
    name=dataonly,
    width=7cm, height=5.5cm,
    xlabel={$x$},
    ylabel={$y$},
    xmin=-0.5, xmax=5.5,
    ymin=0, ymax=5,
    grid=major,
    grid style={UofSC50Black!30},
    axis lines=left,
    xtick={0,1,2,3,4,5},
    ytick={1,2,3,4},
    title={\textbf{Data Points}},
    title style={at={(0.5,1.02)}, font=\bfseries},
    clip=false
]

% Data points
\addplot[only marks, mark=*, mark size=4pt, mark options={fill=UofSCBlack, draw=UofSCBlack}]
    coordinates {(0, 1) (1, 2.5) (2, 1.8) (3, 3.2) (4, 2.9) (5, 4.1)};

% Labels
\node[font=\footnotesize, anchor=south] at (axis cs:0,1.15) {$(x_0, y_0)$};
\node[font=\footnotesize, anchor=south] at (axis cs:5,4.25) {$(x_n, y_n)$};

\end{axis}

% ========== Linear Interpolation ==========
\begin{axis}[
    name=linear,
    at={(dataonly.south east)},
    anchor=south west,
    xshift=1cm,
    width=7cm, height=5.5cm,
    xlabel={$x$},
    ylabel={$y$},
    xmin=-0.5, xmax=5.5,
    ymin=0, ymax=5,
    grid=major,
    grid style={UofSC50Black!30},
    axis lines=left,
    xtick={0,1,2,3,4,5},
    ytick={1,2,3,4},
    title={\textbf{Linear Interpolation}},
    title style={at={(0.5,1.02)}, font=\bfseries},
    clip=false
]

% Linear interpolation (piecewise linear)
\addplot[UofSCGarnet, very thick]
    coordinates {(0, 1) (1, 2.5) (2, 1.8) (3, 3.2) (4, 2.9) (5, 4.1)};

% Data points
\addplot[only marks, mark=*, mark size=4pt, mark options={fill=UofSCBlack, draw=UofSCBlack}]
    coordinates {(0, 1) (1, 2.5) (2, 1.8) (3, 3.2) (4, 2.9) (5, 4.1)};

% Highlight a segment
\draw[UofSCGarnet, thick, dashed] (axis cs:1.5,0) -- (axis cs:1.5,2.15);
\node[circle, fill=UofSCGarnet, minimum size=4pt, inner sep=0pt] at (axis cs:1.5,2.15) {};

\end{axis}

% ========== Polynomial Interpolation ==========
\begin{axis}[
    name=polynomial,
    at={(dataonly.south west)},
    anchor=north west,
    yshift=-1.3cm,
    width=7cm, height=5.5cm,
    xlabel={$x$},
    ylabel={$y$},
    xmin=-0.5, xmax=5.5,
    ymin=0, ymax=5,
    grid=major,
    grid style={UofSC50Black!30},
    axis lines=left,
    xtick={0,1,2,3,4,5},
    ytick={1,2,3,4},
    title={\textbf{Polynomial Interpolation}},
    title style={at={(0.5,1.02)}, font=\bfseries},
    clip=true
]

% Lagrange polynomial through 6 points (degree 5)
% Computed polynomial: p(x) = 1 + 2.2833x - 1.7917x^2 + 0.7583x^3 - 0.1333x^4 + 0.0083x^5
% Simplified approximation for visualization
\addplot[UofSCAtlantic, very thick, domain=0:5, samples=150]
    {1 + 2.2833*x - 1.7917*x^2 + 0.7583*x^3 - 0.1333*x^4 + 0.0083*x^5};

% Data points
\addplot[only marks, mark=*, mark size=4pt, mark options={fill=UofSCBlack, draw=UofSCBlack}]
    coordinates {(0, 1) (1, 2.5) (2, 1.8) (3, 3.2) (4, 2.9) (5, 4.1)};

% Show oscillation label
\node[UofSCAtlantic, font=\footnotesize, align=center] at (axis cs:2.5,0.8) {Single polynomial\\(degree $n-1$)};

\end{axis}

% ========== Cubic Spline Interpolation ==========
\begin{axis}[
    name=spline,
    at={(polynomial.south east)},
    anchor=south west,
    xshift=1cm,
    width=7cm, height=5.5cm,
    xlabel={$x$},
    ylabel={$y$},
    xmin=-0.5, xmax=5.5,
    ymin=0, ymax=5,
    grid=major,
    grid style={UofSC50Black!30},
    axis lines=left,
    xtick={0,1,2,3,4,5},
    ytick={1,2,3,4},
    title={\textbf{Cubic Spline Interpolation}},
    title style={at={(0.5,1.02)}, font=\bfseries},
    clip=true
]

% Natural cubic spline (smooth piecewise cubic)
% Approximate spline using smooth Bezier-like curve
\addplot[UofSCHorseshoe, very thick, smooth, tension=0.55]
    coordinates {(0, 1) (1, 2.5) (2, 1.8) (3, 3.2) (4, 2.9) (5, 4.1)};

% Data points
\addplot[only marks, mark=*, mark size=4pt, mark options={fill=UofSCBlack, draw=UofSCBlack}]
    coordinates {(0, 1) (1, 2.5) (2, 1.8) (3, 3.2) (4, 2.9) (5, 4.1)};

% Show continuity at knot
\draw[UofSCCongaree, thick, ->] (axis cs:2,1.8) -- (axis cs:2.3,1.3);
\node[UofSCCongaree, font=\footnotesize, anchor=north west] at (axis cs:2.3,1.3) {$C^2$ at knots};

\end{axis}

% ========== Comparison inset ==========
\begin{axis}[
    name=comparison,
    at={(spline.south east)},
    anchor=south east,
    xshift=-0.3cm,
    yshift=1.5cm,
    width=4.5cm, height=3.5cm,
    xlabel={$x$},
    xlabel style={font=\footnotesize},
    xmin=0.8, xmax=2.2,
    ymin=1.5, ymax=3,
    grid=major,
    grid style={UofSC50Black!30},
    xtick={1,1.5,2},
    ytick={2,2.5},
    tick label style={font=\tiny},
    title={\footnotesize\textbf{Zoomed: $x \in [1,2]$}},
    title style={at={(0.5,1.0)}, font=\footnotesize\bfseries}
]

% Linear
\addplot[UofSCGarnet, thick]
    coordinates {(1, 2.5) (2, 1.8)};

% Polynomial
\addplot[UofSCAtlantic, thick, domain=1:2, samples=50]
    {1 + 2.2833*x - 1.7917*x^2 + 0.7583*x^3 - 0.1333*x^4 + 0.0083*x^5};

% Spline (approximated)
\addplot[UofSCHorseshoe, thick, smooth, tension=0.55]
    coordinates {(1, 2.5) (1.5, 2.25) (2, 1.8)};

% Data points
\addplot[only marks, mark=*, mark size=2pt, mark options={fill=UofSCBlack}]
    coordinates {(1, 2.5) (2, 1.8)};

\end{axis}

% ========== Method descriptions ==========
\node[anchor=north west, align=left, font=\small] at (0,-8.3) {
    \textcolor{UofSCGarnet}{\textbf{Linear:}} $p(x) = y_i + \dfrac{y_{i+1} - y_i}{x_{i+1} - x_i}(x - x_i)$ for $x \in [x_i, x_{i+1}]$ \quad ($C^0$ continuity)
};
\node[anchor=north west, align=left, font=\small] at (0,-9.1) {
    \textcolor{UofSCAtlantic}{\textbf{Polynomial (Lagrange):}} $p(x) = \displaystyle\sum_{i=0}^{n} y_i \prod_{j \neq i} \frac{x - x_j}{x_i - x_j}$ \quad (Unique degree $n-1$ polynomial)
};
\node[anchor=north west, align=left, font=\small] at (0,-9.9) {
    \textcolor{UofSCHorseshoe}{\textbf{Cubic Spline:}} Piecewise cubic $S_i(x)$ with $S_i(x_i) = y_i$, $S_i(x_{i+1}) = y_{i+1}$, $S'_i, S''_i$ continuous \quad ($C^2$ continuity)
};

\end{tikzpicture}
\caption{Comparison of interpolation methods for discrete data points $(x_i, y_i)$. \textbf{Top left}: The raw data points to be interpolated. \textbf{Top right}: Linear interpolation connects adjacent points with straight line segments; simple but only $C^0$ continuous (discontinuous derivatives at knots). \textbf{Bottom left}: Polynomial interpolation fits a single polynomial of degree $n-1$ through $n$ points; exact at data points but may exhibit oscillation (Runge phenomenon) for high-degree polynomials. \textbf{Bottom right}: Cubic spline interpolation uses piecewise cubic polynomials that are $C^2$ continuous at knot points, providing smooth interpolation without oscillation. The inset compares all three methods in the interval $[1,2]$.}
\label{fig:interpolation}
\end{figure}
